% The PREP operator of the LCU build (paper section 3.1).  Same conventions as
% figures/circuits.tex: wires left to right, a box acts on the wires it spans.
% A filled SQUARE is a uniform control: one angle per value of the controlling
% qubits (a uniformly controlled rotation), not a single control value.
\begin{tikzpicture}[x=1cm,y=1cm]
\tikzset{
  uctrl/.style={rectangle,fill=black,inner sep=0pt,minimum size=1.9mm},
  tnode/.style={draw,rounded corners=1pt,fill=white,inner sep=1.5pt,font=\footnotesize,line width=0.5pt,minimum width=7.5mm},
  tdead/.style={tnode,text=black!45,draw=black!35},
  ang/.style={font=\scriptsize,text=blue!45!black,inner sep=1pt,fill=white},
}
% ═════════════════ (a) the cascade on the index register ═════════════════════
\node[anchor=west,font=\small\bfseries] at (-2.4,3.0)
  {(a) PREP: one uniformly controlled $R_y$ per index qubit};
\foreach \i/\y in {0/1.75,1/1.05,2/0.35,3/-0.35,4/-1.05,5/-1.75}{
  \node[reg] at (-0.15,\y) {$|0\rangle_{\i}$};
  \draw[wire] (0,\y) -- (7.5,\y);}
\node[gate,minimum width=12mm] at (1.3,1.75) {$R_y(\theta^{(0)})$};
\node[uctrl] at (2.7,1.75) {}; \draw[wire] (2.7,1.75) -- (2.7,1.05);
\node[gate,minimum width=12mm] at (2.7,1.05) {$R_y(\theta^{(1)}_i)$};
\node[uctrl] at (4.1,1.75) {}; \node[uctrl] at (4.1,1.05) {}; \draw[wire] (4.1,1.75) -- (4.1,0.35);
\node[gate,minimum width=12mm] at (4.1,0.35) {$R_y(\theta^{(2)}_i)$};
\node[font=\large,text=black!55] at (5.35,-0.35) {$\cdots$};
\foreach \y in {1.75,1.05,0.35,-0.35,-1.05}{\node[uctrl] at (6.5,\y) {};}
\draw[wire] (6.5,1.75) -- (6.5,-1.75);
\node[gate,minimum width=12mm] at (6.5,-1.75) {$R_y(\theta^{(5)}_i)$};
\vsep{2.0}{-2.25}{2.15} \vsep{3.4}{-2.25}{2.15} \vsep{4.8}{-2.25}{2.15} \vsep{7.2}{-2.25}{2.15}
\node[blk] at (1.3,-2.3) {$1$ angle}; \node[blk] at (2.7,-2.3) {$2$};
\node[blk] at (4.1,-2.3) {$4$}; \node[blk] at (6.5,-2.3) {$32$};
\node[st] at (2.0,2.10) {$|\Phi^{(0)}\rangle$}; \node[st] at (3.4,2.10) {$|\Phi^{(1)}\rangle$};
\node[st] at (4.8,2.10) {$|\Phi^{(2)}\rangle$}; \node[st] at (7.2,2.10) {$|\Phi^{(5)}\rangle$};
\node[font=\footnotesize,anchor=north west,text=black!65,align=left,text width=10.4cm] at (-2.4,-2.62)
  {Index register $I$, $m=6$ qubits, $q_0$ the leading bit of $k$. The rotation on
   qubit $j$ is chosen by the $j$ bits above it, so level $j$ carries $2^{j}$
   angles and the whole cascade $2^{m}-1=63$:
   $|\Phi^{(5)}\rangle=\sum_{k}\sqrt{|c_k|/\lVert c\rVert_1}\,|k\rangle$.};
% ═════════════════ (b) the tree that fixes the angles ════════════════════════
\begin{scope}[shift={(8.9,0.15)}]
\node[anchor=west,font=\small\bfseries] at (-0.5,2.85)
  {(b) where the angles come from};
\node[tnode] (r)  at (2.75, 1.65) {$1$};
\node[tnode] (n0) at (1.25, 0.55) {$0.946$};
\node[tnode] (n1) at (4.25, 0.55) {$0.324$};
\node[tnode] (a0) at (0.5,-0.55) {$0.673$};
\node[tnode] (a1) at (2.0,-0.55) {$0.665$};
\node[tnode] (a2) at (3.5,-0.55) {$0.324$};
\node[tdead] (a3) at (5.0,-0.55) {$0$};
\draw[wire] (r) -- (n0); \draw[wire] (r) -- (n1);
\draw[wire] (n0) -- (a0); \draw[wire] (n0) -- (a1);
\draw[wire] (n1) -- (a2); \draw[wire,black!35] (n1) -- (a3);
\node[ang] at (2.75,1.10) {$\theta^{(0)}=37.8^{\circ}$};
\node[ang] at (1.25,0.00) {$89.3^{\circ}$};
\node[ang] at (4.25,0.00) {$0^{\circ}$};
\node[ang] at (0.5,-1.05) {$71.8^{\circ}$};
\node[ang] at (2.0,-1.05) {$73.6^{\circ}$};
\node[ang] at (3.5,-1.05) {$79.8^{\circ}$};
\node[font=\scriptsize,anchor=north] at (2.0,-1.35) {$\cdots$ three more levels $\cdots$};
\node[font=\scriptsize,anchor=north,text=black!55] at (5.0,-1.05) {$k>40$};
\node[font=\scriptsize,anchor=north west,text=black!75,align=left,text width=6.3cm] at (-0.5,-2.05)
  {A node carries the norm of the coefficients below it,
   $A(v)=\sqrt{A(v_{\rm L})^{2}+A(v_{\rm R})^{2}}$, the leaves carry
   $\pm\sqrt{|c_k|/\lVert c\rVert_1}$, and the rotation that splits a node is
   $\theta_v=2\arctan\big(A(v_{\rm R})/A(v_{\rm L})\big)$ taken in the full circle.
   The $63$ angles are exactly enough to carry the $63$ relative signs too, so the
   $15$ negative $c_k$ need no diagonal of their own; $23$ of the $64$ index states
   have no coefficient, and the branches under them are dropped.};
\end{scope}
% ═════════════════ the two notes along the foot ══════════════════════════════
\begin{scope}[shift={(0,-5.15)}]
  \node[anchor=west,font=\small\bfseries] at (-2.4,0.72) {(c) the system register, for comparison};
  \node[reg] at (-0.15,0) {$|0\rangle^{\otimes n}_{S}$};
  \draw[wire] (0,0) -- (3.0,0); \slashwire{0.3}{0}{$n$}
  \node[gate,minimum width=10mm] at (1.5,0) {$H^{\otimes n}$};
  \node[font=\small,anchor=west] at (3.15,0) {$=|u\rangle$};
  \node[font=\footnotesize,anchor=north west,text=black!65,align=left,text width=10.4cm] at (-2.4,-0.62)
    {$\mathbf{b}=\delta\mathbf{1}$ is uniform, so the system register's own preparation is
     exact and costs one layer of $n$ Hadamards: no arbitrary state is loaded anywhere
     outside the index register.};
\end{scope}
\end{tikzpicture}
